\section{Distribución Audiovisual}
\label{sec:distribucion_av}

Una vez generada la señal de programa por el mezclador, esta debe distribuirse hasta
la audiencia. Esta etapa difiere fundamentalmente de la contribución: mientras que en
la contribución se prioriza la calidad y la latencia mínima para preservar la señal
sin degradación, en la distribución el objetivo es llegar al máximo número de
espectadores con los recursos de red disponibles, lo que obliga a comprimir la señal
de forma agresiva y a equilibrar el compromiso entre latencia, calidad y escalabilidad.

\subsection{Interfaces y protocolos de distribución}

\begin{itemize}
  \item \textbf{RTMP (\textit{Real-Time Messaging Protocol}):} protocolo TCP de Adobe,
  estándar de ingesta para plataformas como YouTube Live o Twitch con latencia de
  2--5~s. Aunque obsoleto en reproducción desde la desaparición de Flash, sigue siendo
  el protocolo de ingesta más soportado, incluyendo Voctomix mediante FFmpeg, con una
  adopción profesional del 58\% en 2025~\cite{haivision_broadcast_report_2025}.

  \item \textbf{SRT (\textit{Secure Reliable Transport}):} protocolo de código abierto
  de Haivision con latencia inferior a 1~s, corrección de errores ARQ y cifrado
  AES-128/256, diseñado para transmisión sobre redes no
  fiables~\cite{sharabayko2024_srt}. Es actualmente el protocolo de transporte en
  directo más utilizado entre profesionales, con una adopción del \textbf{77\%} en
  2025, superando a RTMP~\cite{haivision_broadcast_report_2025}.

  \item \textbf{HLS (\textit{HTTP Live Streaming}):} protocolo de Apple que divide el
  vídeo en segmentos servidos mediante HTTP~\cite{rfc8216}. Latencia de 10--30~s en
  modo estándar, reducible a 2--5~s con \textit{Low-Latency HLS}. Es el protocolo de
  reproducción más compatible, con soporte nativo en prácticamente todos los
  dispositivos y CDN.

  \item \textbf{MPEG-DASH (\textit{Dynamic Adaptive Streaming over HTTP}):} estándar
  abierto ISO/IEC~23009-1~\cite{iso_mpegdash} equivalente a HLS pero sin dependencia
  de proveedor. Divide el contenido en segmentos HTTP con ajuste adaptativo de calidad
  (ABR), y es empleado por plataformas como YouTube, Netflix y Amazon Prime. Puede
  generarse con FFmpeg como alternativa abierta a HLS en el pipeline de Voctomix.

  \item \textbf{WebRTC (\textit{Web Real-Time Communication}):} estándar W3C para
  comunicación multimedia en tiempo real desde el navegador, con latencias de
  100--500~ms. Su escalabilidad para audiencias masivas es limitada, pero resulta ideal
  para monitorización remota del sistema de producción y retornos entre el equipo
  técnico.
\end{itemize}

\begin{table}[H]
  \centering
  \caption{Comparativa de protocolos de distribución audiovisual en directo.}
  \label{tab:distribucion_protocolos}
  \begin{tabular}{|l|c|l|c|l|}
    \hline
    \textbf{Protocolo} & \textbf{Latencia} & \textbf{Transporte} & \textbf{Cifrado} & \textbf{Uso principal} \\
    \hline
    RTMP      & 2--5~s      & TCP    & No (RTMPS sí) & Ingesta a plataformas \\
    SRT       & $<$1~s      & UDP    & AES-128/256   & Contribución/distribución IP \\
    HLS       & 2--30~s     & HTTP   & HTTPS         & Reproducción masiva (CDN) \\
    MPEG-DASH & 2--30~s     & HTTP   & HTTPS         & Streaming adaptativo abierto \\
    WebRTC    & 100--500~ms & UDP    & DTLS/SRTP     & Monitorización, retornos \\
    \hline
  \end{tabular}
\end{table}

\subsection{Códecs de Vídeo}

En la etapa de distribución se emplean códecs de alta eficiencia que comprimen la señal
para reducir el ancho de banda necesario hasta el espectador final, asumiendo cierta
pérdida de calidad a cambio de viabilizar el transporte:

\begin{itemize}
  \item \textbf{H.264 / AVC} (\textit{Advanced Video Coding}, ITU-T H.264):
  el códec más extendido en la actualidad y la línea de base de la industria. Ofrece
  buena calidad a tasas de bits moderadas y es soportado por la práctica totalidad de
  dispositivos, plataformas y hardware de codificación desde 2003. Voctomix lo emplea
  como códec de salida por defecto mediante FFmpeg~\cite{itu_h264}.

  \item \textbf{H.265 / HEVC} (\textit{High Efficiency Video Coding}): sucesor de
  H.264 que aproximadamente \textbf{duplica la eficiencia de compresión}, permitiendo
  la misma calidad perceptual a la mitad del \textit{bitrate} (tasa de bits)~\cite{itu_h265}. Su codificación
  es 5--8 veces más lenta que H.264, lo que lo hace adecuado para streaming en tiempo
  real cuando se dispone de aceleración hardware (GPU o ASIC). Los dispositivos
  fabricados tras 2015 disponen en general de soporte de decodificación por hardware.

  \item \textbf{AV1}: códec abierto desarrollado por la Alliance for Open Media que
  mejora la eficiencia de H.265 en aproximadamente un \textbf{50\%} adicional~\cite{chen2024_codec_comparison}.
  Su limitación en producción en directo es la velocidad de codificación: AV1 tarda entre
  5 y 10 veces más que H.265 en codificar el mismo contenido, lo que restringe su uso
  en tiempo real a plataformas con aceleración hardware dedicada. YouTube, Netflix y
  Facebook ya codifican gran parte de su contenido en AV1 para streaming bajo demanda.

  \item \textbf{AAC / Opus}: códecs de audio que acompañan habitual\-mente a los
  anteriores en la distribución. AAC (\textit{Advanced Audio Coding}) es el estándar
  en distribución broadcast y plataformas de streaming; Opus es el códec preferido en
  aplicaciones WebRTC por su eficiencia a bajas tasas de bits.
\end{itemize}

\begin{table}[H]
  \centering
  \caption{Comparativa de códecs de vídeo para distribución en directo.}
  \label{tab:codecs_distribucion}
  \begin{tabular}{|l|c|c|c|l|}
    \hline
    \textbf{Códec} & \textbf{Eficiencia vs H.264} & \textbf{Vel. encoding} & \textbf{Compatibilidad} & \textbf{Streaming RT} \\
    \hline
    H.264/AVC & Base       & Muy rápido  & Universal    & Sí \\
    H.265/HEVC & $\times$2 & 5--8x más lento & Alta (post 2015) & Sí (con HW) \\
    AV1        & $\times$3 & 5--10x más lento & Media (creciendo) & Solo con HW \\
    \hline
  \end{tabular}
\end{table}

\subsection{Formatos}

El formato contenedor de distribución define cómo se encapsulan los flujos de vídeo
y audio codificados para su transporte hasta el reproductor final:

\begin{itemize}
  \item \textbf{MPEG-TS (\textit{Transport Stream}, ISO 13818-1):} estándar diseñado
  específicamente para transmisión broadcast con tolerancia a errores. Es el contenedor
  asociado a HLS, la emisión por satélite y la TDT.

  \item \textbf{MP4 / ISOBMFF (\textit{ISO Base Media File Format}):} contenedor de
  referencia para distribución por HTTP progresivo y \textit{streaming} adaptativo
  (MPEG-DASH, HLS fragmentado). Es el formato de almacenamiento estándar para
  grabaciones y distribución bajo demanda.

  \item \textbf{FLV (\textit{Flash Video}):} contenedor legado asociado a RTMP.
  Aunque técnicamente obsoleto como formato de reproducción desde la retirada de Adobe
  Flash en 2020, sigue siendo el contenedor utilizado internamente en el protocolo RTMP
  para la ingesta hacia plataformas de streaming en directo.
\end{itemize}
